\documentclass[12pt,a4paper]{amsart}

\usepackage{amsmath,amssymb,amsthm}
\usepackage{mathtools}
\usepackage{hyperref}
\usepackage{geometry}
\usepackage[ngerman]{babel}
\geometry{margin=2.5cm}

% -------------------------------------------------------
\newtheorem{theorem}{Satz}[section]
\newtheorem{lemma}[theorem]{Lemma}
\newtheorem{corollary}[theorem]{Korollar}
\newtheorem{proposition}[theorem]{Proposition}
\theoremstyle{definition}
\newtheorem{definition}[theorem]{Definition}
\newtheorem{remark}[theorem]{Bemerkung}

\newcommand{\Z}{\mathbb{Z}}

% -------------------------------------------------------
\begin{document}
% -------------------------------------------------------

\title{Kein Semiprimes ist ein Giuga-Pseudoprim:\\ Ein elementares Ordnungsargument}

\author{Kurt Ingwer}
\address{Specialist Mathematics Project, Version Build 44}
\email{specialist-maths@localhost}
\date{\today}

\begin{abstract}
Ein \emph{Giuga-Pseudoprim} ist eine zusammengesetzte ganze Zahl, die Giugas
Primalitätskriterium von 1950 erfüllt. Ein \emph{Semiprim} ist ein Produkt genau zweier
(nicht notwendig verschiedener) Primzahlen. Wir beweisen, dass kein Semiprim
$n = p \cdot q$ mit $p < q$ prim ein Giuga-Pseudoprim ist. Das Argument ist eine
einfache Ordnungsungleichung: Die schwache Giuga-Bedingung für die größere Primzahl
erzwingt $q \mid (p-1)$, aber $p - 1 < q$, also $p - 1 = 0$ und $p = 1$, was der
Primalität widerspricht. Zusammen mit der Quadratfreiheit von Giuga-Pseudoprimen folgt,
dass jedes Giuga-Pseudoprim mindestens \emph{drei} verschiedene Primfaktoren haben muss.
\end{abstract}

\maketitle

% -------------------------------------------------------
\section{Einleitung}
% -------------------------------------------------------

Die Giuga-Bedingungen für eine zusammengesetzte Zahl~$n$ bezüglich eines Primteilers
$p \mid n$ sind:
\begin{enumerate}
  \item[\textup{(S)}] $p \mid \tfrac{n}{p} - 1$ \quad (schwache Bedingung),
  \item[\textup{(St)}] $(p-1) \mid \tfrac{n}{p} - 1$ \quad (starke Bedingung).
\end{enumerate}
Eine zusammengesetzte Zahl~$n$, die beide Bedingungen für jeden Primfaktor erfüllt,
heißt \emph{Giuga-Pseudoprim} \cite{Borwein1996}.

Es ist bereits bekannt (und elementar), dass jedes Giuga-Pseudoprim quadratfrei ist
\cite{Giuga1950}. Die vorliegende Notiz schließt zusammengesetzte Zahlen mit genau
zwei Primfaktoren aus.

% -------------------------------------------------------
\section{Hauptresultat}
% -------------------------------------------------------

\begin{theorem}
\label{thm:keinpq}
Es existiert kein Giuga-Pseudoprim der Form $n = p \cdot q$ mit $p < q$ prim.
\end{theorem}

\begin{proof}
Angenommen, $n = pq$ ist ein Giuga-Pseudoprim mit Primzahlen $p < q$. Die Quotienten
sind $\tfrac{n}{p} = q$ und $\tfrac{n}{q} = p$.

Wir wenden die schwache Bedingung~(S) auf beide Primfaktoren an:
\begin{align}
  p &\mid q - 1, \label{eq:Sp} \\
  q &\mid p - 1. \label{eq:Sq}
\end{align}

Aus~\eqref{eq:Sq}: Da $p < q$, gilt $0 \le p - 1 < q$.
Das einzige nichtnegative Vielfache von~$q$, das kleiner als~$q$ ist, ist~$0$.
Daher erzwingt $q \mid (p - 1)$ die Konsequenz $p - 1 = 0$, also $p = 1$.
Aber $1$ ist keine Primzahl — Widerspruch.
\end{proof}

\begin{remark}
Das Argument verwendet nur die schwache Bedingung~(S) für die \emph{größere}
Primzahl~$q$. Bedingung~(S) für~$p$ (Gleichung~\eqref{eq:Sp}) wird für den
Widerspruch nicht benötigt.
\end{remark}

\begin{remark}
Der Beweis ist ein reines Ordnungsargument: Die Teilbarkeit $q \mid (p-1)$ zusammen
mit der strikten Ungleichung $p - 1 < q$ erzwingt $p - 1 = 0$. Jenseits grundlegender
Teilbarkeitseigenschaften wird keine Arithmetik benötigt.
\end{remark}

% -------------------------------------------------------
\section{Folgerung: Mindestens drei Primfaktoren}
% -------------------------------------------------------

\begin{corollary}
\label{cor:drei}
Jedes Giuga-Pseudoprim hat mindestens drei verschiedene Primfaktoren.
\end{corollary}

\begin{proof}
Ein Giuga-Pseudoprim ist quadratfrei \cite{Ingwer2026qfrei}, also ein Produkt
verschiedener Primzahlen. Mit einem Primfaktor wäre $n$ prim — nicht zusammengesetzt.
Mit zwei Primfaktoren greift Satz~\ref{thm:keinpq}. Also werden mindestens drei benötigt.
\end{proof}

\begin{remark}
Die untere Schranke von drei ist insofern scharf, als die Bedingungen drei Faktoren
nicht sofort ausschließen. Ein weiteres Argument zeigt jedoch, dass auch drei Faktoren
unmöglich sind \cite{Ingwer2026a}, sodass tatsächlich mindestens \emph{vier}
Primfaktoren erforderlich sind. Die aktuelle beste Schranke, von Bednarek und Borwein
\cite{Bednarek2014}, liegt bei $59$ Primfaktoren.
\end{remark}

% -------------------------------------------------------
\section{Symmetriebeobachtung}
% -------------------------------------------------------

Falls $n = pq$ mit $p = q$ (ein Primzahlquadrat), ist $n$ nicht quadratfrei und
bereits ausgeschlossen. Für $p < q$ gilt obiges Argument. Es besteht damit eine
angenehme Symmetrie: Der Zweiprimensfall wird durch eine Kombination eliminiert:
\begin{itemize}
  \item Quadratfreiheit (schließt $p = q$ aus), und
  \item Das Ordnungsargument (schließt $p < q$ aus).
\end{itemize}

Dieser zweistufige Ausschluss ist charakteristisch für die Giuga-Theorie: Strukturelle
Eigenschaften (Quadratfreiheit) verbinden sich mit Ordnungsungleichungen, um Widersprüche
zu erzeugen.

% -------------------------------------------------------
\begin{thebibliography}{10}

\bibitem{Giuga1950}
G.~Giuga,
\emph{Su una presumibile propriet\`a caratteristica dei numeri primi},
Ist.\ Lombardo Sci.\ Lett.\ Rend.\ Cl.\ Sci.\ Mat.\ Nat.\ \textbf{83}
(1950), 511--528.

\bibitem{Borwein1996}
J.~M.~Borwein, P.~B.~Borwein, R.~Girgensohn und S.~Parnes,
\emph{Giuga's conjecture on primality},
Amer.\ Math.\ Monthly \textbf{103} (1996), Nr.\ 1, 40--50.

\bibitem{Bednarek2014}
M.~Bednarek,
\emph{On the size of Giuga pseudoprimes},
Preprint (2014).

\bibitem{Ingwer2026qfrei}
K.~Ingwer,
\emph{Jedes Giuga-Pseudoprim ist quadratfrei},
Specialist Mathematics Project (2026).

\bibitem{Ingwer2026a}
K.~Ingwer,
\emph{Keine zusammengesetzte Zahl mit genau drei Primfaktoren ist ein Giuga-Pseudoprim},
Specialist Mathematics Project (2026).

\end{thebibliography}

\end{document}
